\documentclass[12pt, a4paper]{article}

% -------------------------------------------------------
% PACKAGES
% -------------------------------------------------------
\usepackage[a4paper, margin=1in, headheight=15pt]{geometry}
\usepackage{fontenc}
\usepackage{inputenc}
\usepackage[T1]{fontenc}
\usepackage{lmodern}
\usepackage{microtype}
\usepackage{graphicx}
\usepackage{xcolor}
\usepackage{listings}
\usepackage{booktabs}
\usepackage{array}
\usepackage{enumitem}
\usepackage{titlesec}
\usepackage{fancyhdr}
\usepackage{tcolorbox}
\usepackage{tikz}
\usepackage{hyperref}
\usepackage{setspace}
\usepackage{parskip}
\usepackage{amsmath}
\usepackage{mdframed}
\usepackage{caption}

\tcbuselibrary{skins, breakable, listings}

% -------------------------------------------------------
% COLORS
% -------------------------------------------------------
\definecolor{primaryblue}{RGB}{30, 90, 160}
\definecolor{secondaryblue}{RGB}{70, 130, 200}
\definecolor{lightblue}{RGB}{220, 235, 252}
\definecolor{codered}{RGB}{180, 30, 30}
\definecolor{codegreen}{RGB}{30, 130, 30}
\definecolor{codegray}{RGB}{100, 100, 100}
\definecolor{codebg}{RGB}{245, 247, 250}
\definecolor{keywordcolor}{RGB}{0, 80, 180}
\definecolor{stringcolor}{RGB}{163, 21, 21}
\definecolor{commentcolor}{RGB}{80, 120, 80}
\definecolor{warnbox}{RGB}{255, 243, 220}
\definecolor{warnborder}{RGB}{210, 140, 30}
\definecolor{successbox}{RGB}{220, 245, 220}
\definecolor{successborder}{RGB}{50, 150, 50}
\definecolor{sectionbg}{RGB}{240, 245, 255}

% -------------------------------------------------------
% HYPERREF SETUP
% -------------------------------------------------------
\hypersetup{
    colorlinks=true,
    linkcolor=primaryblue,
    urlcolor=secondaryblue,
    pdftitle={Factory Method Design Pattern},
    pdfauthor={Smart Transportation Management System Report},
}

% -------------------------------------------------------
% PAGE STYLE
% -------------------------------------------------------
\pagestyle{fancy}
\fancyhf{}
\fancyhead[L]{\small\textcolor{primaryblue}{\textbf{Factory Method Design Pattern}}}
\fancyhead[R]{\small\textcolor{codegray}{Smart Transportation Management System}}
\fancyfoot[C]{\textcolor{codegray}{\thepage}}
\renewcommand{\headrulewidth}{0.5pt}
\renewcommand{\headrule}{\hbox to\headwidth{\color{primaryblue}\leaders\hrule height \headrulewidth\hfill}}

% -------------------------------------------------------
% SECTION FORMATTING
% -------------------------------------------------------
\titleformat{\section}
  {\Large\bfseries\color{primaryblue}}
  {\thesection}{1em}{}[{\color{primaryblue}\titlerule[1pt]}]

\titleformat{\subsection}
  {\large\bfseries\color{secondaryblue}}
  {\thesubsection}{1em}{}

\titleformat{\subsubsection}
  {\normalsize\bfseries\color{primaryblue}}
  {\thesubsubsection}{1em}{}

\titlespacing*{\section}{0pt}{18pt}{8pt}
\titlespacing*{\subsection}{0pt}{12pt}{5pt}
\titlespacing*{\subsubsection}{0pt}{10pt}{4pt}

% -------------------------------------------------------
% JAVA LISTINGS STYLE
% -------------------------------------------------------
\lstdefinestyle{javaStyle}{
    language=Java,
    backgroundcolor=\color{codebg},
    basicstyle=\ttfamily\small,
    keywordstyle=\bfseries\color{keywordcolor},
    stringstyle=\color{stringcolor},
    commentstyle=\itshape\color{commentcolor},
    numberstyle=\tiny\color{codegray},
    numbers=left,
    numbersep=8pt,
    stepnumber=1,
    breaklines=true,
    breakatwhitespace=false,
    showstringspaces=false,
    frame=single,
    framesep=4pt,
    rulecolor=\color{secondaryblue},
    tabsize=4,
    captionpos=b,
    xleftmargin=15pt,
    xrightmargin=5pt,
    morekeywords={interface, implements, abstract, override, extends},
}

\lstset{style=javaStyle}

% -------------------------------------------------------
% CUSTOM BOXES
% -------------------------------------------------------
\newtcolorbox{definitionbox}[1]{
    colback=lightblue,
    colframe=primaryblue,
    fonttitle=\bfseries\color{white},
    title=#1,
    arc=4pt,
    boxrule=1pt,
    left=8pt, right=8pt, top=6pt, bottom=6pt,
}

\newtcolorbox{warningbox}{
    colback=warnbox,
    colframe=warnborder,
    fonttitle=\bfseries,
    title={\textcolor{warnborder}{Before Refactoring --- Problem}}  ,
    arc=4pt,
    boxrule=1pt,
    left=8pt, right=8pt, top=6pt, bottom=6pt,
}

\newtcolorbox{successboxenv}{
    colback=successbox,
    colframe=successborder,
    fonttitle=\bfseries,
    title={\textcolor{successborder}{After Refactoring --- Solution}},
    arc=4pt,
    boxrule=1pt,
    left=8pt, right=8pt, top=6pt, bottom=6pt,
}

\newtcolorbox{infobox}[1]{
    colback=sectionbg,
    colframe=secondaryblue,
    fonttitle=\bfseries\color{primaryblue},
    title=#1,
    arc=4pt,
    boxrule=0.8pt,
    left=8pt, right=8pt, top=6pt, bottom=6pt,
}

% -------------------------------------------------------
% DOCUMENT
% -------------------------------------------------------
\begin{document}

% ========================
% TITLE PAGE
% ========================
\begin{titlepage}
    \centering
    \vspace*{1.5cm}

    {\Huge\bfseries\color{primaryblue} Factory Method\\[6pt] Design Pattern}

    \vspace{0.4cm}
    {\large\color{codegray}\rule{10cm}{1pt}}

    \vspace{0.6cm}
    {\Large\color{secondaryblue} Theory, Application \& Refactoring Report}

    \vspace{0.4cm}
    {\normalsize\color{codegray} Smart Transportation Management System --- Java}

    \vspace{2.5cm}

    \begin{tcolorbox}[
        colback=lightblue,
        colframe=primaryblue,
        width=10cm,
        arc=6pt,
        boxrule=1pt,
        halign=center
    ]
        \centering
        \textbf{\color{primaryblue}Pattern Category:} Creational Design Pattern\\[4pt]
        \textbf{\color{primaryblue}Language:} Java\\[4pt]
        \textbf{\color{primaryblue}Scope:} Object Creation \& Factory Abstraction
    \end{tcolorbox}

    \vfill

    {\small\color{codegray} Prepared for: Software Design Patterns Course}\\[4pt]
    {\small\color{codegray} \today}
\end{titlepage}

% ========================
% TABLE OF CONTENTS
% ========================
\tableofcontents
\newpage

% ========================
% PART 1: THEORY
% ========================
\section{Introduction to Design Patterns}

Before understanding the Factory Method Pattern specifically, it is important to understand what design patterns are and why software engineers use them.

\begin{definitionbox}{What is a Design Pattern?}
A \textbf{design pattern} is a reusable, well-tested solution to a commonly occurring problem in software design. It is not finished code you can paste directly into your program --- rather, it is a \textit{blueprint} or \textit{template} that describes how to solve a problem in a structured way.

Design patterns were popularised by the ``Gang of Four'' (GoF) in their 1994 book \textit{Design Patterns: Elements of Reusable Object-Oriented Software}.
\end{definitionbox}

\vspace{6pt}

Design patterns are grouped into three main categories:

\begin{center}
\begin{tabular}{>{\bfseries\color{primaryblue}}l p{9cm}}
\toprule
\textbf{\color{black}Category} & \textbf{Description} \\
\midrule
Creational & Deal with \textit{object creation}. Control how and when objects are instantiated. \\[4pt]
Structural & Deal with \textit{object composition}. Define how objects and classes are assembled. \\[4pt]
Behavioural & Deal with \textit{communication} between objects and responsibilities. \\
\bottomrule
\end{tabular}
\end{center}

\vspace{6pt}
The \textbf{Factory Method Pattern} belongs to the \textbf{Creational} category.

% ========================
\section{Factory Method Design Pattern --- Theory}

\subsection{What Is the Factory Method Pattern?}

\begin{definitionbox}{Definition}
The \textbf{Factory Method Pattern} defines an interface (or abstract class) for creating an object, but lets \textit{subclasses} decide which class to instantiate. The factory method defers object instantiation to subclasses.
\end{definitionbox}

\vspace{6pt}

In simple terms: instead of writing \texttt{new SomeClass()} directly inside your main or client code, you call a \textit{factory method} that decides which object to create and returns it. The client code never needs to know the details of \textit{how} the object was built.

Think of it like ordering food at a restaurant. You say ``I want a burger'', and the kitchen (factory) prepares it for you. You do not need to know the recipe, ingredients, or cooking method. You simply receive the finished product.

\subsection{Why Do We Need It?}

Consider this typical beginner code without any pattern:

\begin{lstlisting}[caption={Direct object creation --- tightly coupled code}]
// In main() or client code:
if (userChoice.equals("Bus")) {
    Vehicle v = new Bus();       // tightly coupled
} else if (userChoice.equals("Taxi")) {
    Vehicle v = new Taxi();      // tightly coupled
} else if (userChoice.equals("Scooter")) {
    Vehicle v = new ElectricScooter(); // tightly coupled
}
\end{lstlisting}

\vspace{4pt}

\begin{warningbox}
\textbf{Problems with the above approach:}
\begin{itemize}[noitemsep]
    \item Every time a new vehicle type is added (e.g., Train, Boat), the \texttt{main()} class must be modified.
    \item Object creation logic is mixed with business logic --- hard to read and maintain.
    \item Unit testing is difficult because \texttt{main()} depends on concrete classes.
    \item Violates the \textbf{Open/Closed Principle}: classes should be \textit{open for extension} but \textit{closed for modification}.
\end{itemize}
\end{warningbox}

\subsection{How Does the Factory Method Pattern Solve This?}

\begin{successboxenv}
\textbf{With the Factory Method Pattern:}
\begin{itemize}[noitemsep]
    \item A dedicated factory class handles object creation.
    \item The client code only calls \texttt{factory.createVehicle()} --- it never uses \texttt{new} directly.
    \item Adding a new vehicle means adding a new factory class \textit{only} --- existing code is untouched.
    \item Each factory class has a single responsibility: create one type of object.
\end{itemize}
\end{successboxenv}

\subsection{Key Benefits}

\begin{enumerate}[leftmargin=*, label=\textbf{\arabic*.}]
    \item \textbf{Loose Coupling} \\
    The client code depends on the abstract \texttt{VehicleFactory} and \texttt{Vehicle} interface, not on any specific class like \texttt{Bus} or \texttt{Taxi}. Changing the internal implementation of \texttt{Bus} has zero impact on the client.

    \item \textbf{Open/Closed Principle (OCP)} \\
    The system is open for extension (add \texttt{TrainFactory}) but closed for modification (you do not touch existing factories or \texttt{main()}).

    \item \textbf{Single Responsibility Principle (SRP)} \\
    Each factory class has exactly one job: creating one type of vehicle. This makes the code easier to understand, test, and maintain.

    \item \textbf{Scalability} \\
    As the number of vehicle types grows (5, 10, 50), the architecture stays clean. You just keep adding factory classes without restructuring anything.

    \item \textbf{Testability} \\
    Because factories are separate classes, you can mock or replace them in unit tests without touching production code.

    \item \textbf{Readability} \\
    The intent of the code becomes clear. \texttt{new BusFactory().createVehicle()} reads like plain English: ``use the bus factory to create a vehicle.''
\end{enumerate}

\subsection{When Should You Use the Factory Method Pattern?}

Use this pattern when:

\begin{itemize}[leftmargin=*, noitemsep]
    \item You do not know ahead of time which class you will need to instantiate.
    \item You want subclasses to control which objects are created.
    \item Your code creates objects of the same family (vehicles, shapes, users) but the exact type varies.
    \item You want to centralise and encapsulate creation logic.
    \item You anticipate that the number of concrete types will grow over time.
\end{itemize}

\subsection{Structure of the Factory Method Pattern}

The pattern consists of four roles:

\begin{center}
\begin{tabular}{>{\bfseries\color{primaryblue}}l l p{7cm}}
\toprule
\textbf{\color{black}Role} & \textbf{Also Called} & \textbf{Responsibility} \\
\midrule
Product & Interface & Defines the common interface all products must follow \\[4pt]
Concrete Product & Implementation & The real object that is created (Bus, Taxi, etc.) \\[4pt]
Creator & Abstract Factory & Declares the factory method; may also define shared logic \\[4pt]
Concrete Creator & Concrete Factory & Overrides the factory method to return a specific product \\
\bottomrule
\end{tabular}
\end{center}

\subsection{How to Implement the Factory Method Pattern --- Step by Step}

\begin{infobox}{Implementation Guide}
\textbf{Step 1:} Define a \textbf{Product Interface} that all products will implement.\\[4pt]
\textbf{Step 2:} Create \textbf{Concrete Product} classes that implement the interface.\\[4pt]
\textbf{Step 3:} Create an \textbf{Abstract Creator} (abstract class) with an abstract \texttt{createProduct()} factory method.\\[4pt]
\textbf{Step 4:} Create \textbf{Concrete Creator} subclasses, each overriding \texttt{createProduct()} to return one specific product.\\[4pt]
\textbf{Step 5:} In the \textbf{Client Code}, use the abstract creator type and call the factory method. Never use \texttt{new ConcreteProduct()} directly.
\end{infobox}

\newpage

% ========================
% PART 2: CHANGES
% ========================
\section{Refactoring the Smart Transportation Management System}

\subsection{Overview of Changes}

The original code provided a working system with a \texttt{Vehicle} interface and four concrete vehicle classes. However, it lacked any structural design and left object creation entirely unmanaged. The following table summarises all the changes made during refactoring.

\begin{center}
\begin{tabular}{>{\bfseries}l l l p{5.5cm}}
\toprule
\textbf{Component} & \textbf{Original} & \textbf{After Refactoring} & \textbf{Reason} \\
\midrule
\texttt{Vehicle} interface & Present & Kept unchanged & Already well-defined \\[4pt]
\texttt{Bus}, \texttt{Taxi}, etc. & Present & Minor improvements & Core logic was correct \\[4pt]
\texttt{VehicleFactory} & Absent & Added (abstract class) & Core of the pattern \\[4pt]
\texttt{BusFactory} & Absent & Added & Creates \texttt{Bus} only \\[4pt]
\texttt{TaxiFactory} & Absent & Added & Creates \texttt{Taxi} only \\[4pt]
\texttt{MotorcycleDeliveryFactory} & Absent & Added & Creates \texttt{MotorcycleDelivery} \\[4pt]
\texttt{ElectricScooterFactory} & Absent & Added & Creates \texttt{ElectricScooter} \\[4pt]
\texttt{Main} / Client Code & Absent & Added with menu & Demonstrates pattern usage \\
\bottomrule
\end{tabular}
\end{center}

\subsection{Change 1 --- Vehicle Interface Kept Intact}

\subsubsection{Original Code}

\begin{lstlisting}[caption={Original Vehicle interface}]
interface Vehicle {
    void startTrip();
    double calculateFare(double distance);
    void assignRoute(String route);
    void getVehicleInfo();
    String getVehicleType();
}
\end{lstlisting}

\subsubsection{Decision}

The \texttt{Vehicle} interface was \textbf{not modified}. It already correctly defines the contract that all vehicle types must follow. This is precisely the \textbf{Product Interface} role in the Factory Method Pattern. Keeping it unchanged confirms that the original design had a good foundation --- it only lacked the creation layer on top.

\subsection{Change 2 --- Concrete Vehicle Classes Improved}

\subsubsection{Original Code (example: Bus)}

\begin{lstlisting}[caption={Original Bus class --- getVehicleInfo()}]
@Override
public void getVehicleInfo() {
    System.out.println("Vehicle: Bus");
    System.out.println("Capacity: " + capacity);
    System.out.println("Fare per KM: " + farePerKm);
}
\end{lstlisting}

\subsubsection{Refactored Code}

\begin{lstlisting}[caption={Refactored Bus class --- improved output formatting}]
@Override
public void getVehicleInfo() {
    System.out.println("Vehicle     : Bus");
    System.out.println("Capacity    : " + capacity + " passengers");
    System.out.println("Fare per KM : " + farePerKm + " BDT");
}
\end{lstlisting}

\subsubsection{Why This Change?}

The output formatting was improved for two reasons. First, aligned labels (using spaces) make the printed information easier to read at a glance, which matters when the system displays multiple vehicle types in sequence. Second, adding units (\texttt{passengers}, \texttt{BDT}) makes the output self-explanatory --- a reader does not need to guess whether capacity is in seats, tonnes, or cubic metres, or whether the fare is in Taka, Dollars, or any other currency.

The core logic (\texttt{farePerKm}, \texttt{capacity} fields, method implementations) was kept identical, as it was already correct.

\subsection{Change 3 --- Abstract Creator: VehicleFactory (New Addition)}

\subsubsection{Code Added}

\begin{lstlisting}[caption={Abstract VehicleFactory --- the heart of the pattern}]
abstract class VehicleFactory {

    public abstract Vehicle createVehicle();

    public Vehicle orderVehicle(String route) {
        Vehicle vehicle = createVehicle();
        System.out.println("\n[Factory] Creating a "
            + vehicle.getVehicleType() + "...");
        vehicle.assignRoute(route);
        return vehicle;
    }
}
\end{lstlisting}

\subsubsection{Why This Change?}

This is the most critical addition. \texttt{VehicleFactory} is the \textbf{Abstract Creator} in the Factory Method Pattern. Here is why each part matters:

\begin{itemize}[leftmargin=*, label=\textcolor{primaryblue}{\textbullet}]
    \item \textbf{\texttt{public abstract Vehicle createVehicle()}} --- This is the factory method itself. It is declared \texttt{abstract} so that each subclass \textit{must} provide its own implementation. The return type is \texttt{Vehicle} (the interface), not a concrete class --- this ensures the client code remains decoupled from any specific vehicle type.

    \item \textbf{\texttt{public Vehicle orderVehicle(String route)}} --- This is a convenience method that uses the factory method. It calls \texttt{createVehicle()}, which triggers whichever subclass overrides it, and then automatically assigns a route. This demonstrates the pattern's power: the shared logic lives in the abstract class, while object creation is delegated to subclasses.
\end{itemize}

Without this class, there is no factory --- the system is just a collection of unrelated vehicle classes with no controlled creation strategy.

\subsection{Change 4 --- Concrete Creators: Four Factory Classes (New Additions)}

\subsubsection{Code Added}

\begin{lstlisting}[caption={Four concrete factory classes}]
class BusFactory extends VehicleFactory {
    @Override
    public Vehicle createVehicle() {
        return new Bus();
    }
}

class TaxiFactory extends VehicleFactory {
    @Override
    public Vehicle createVehicle() {
        return new Taxi();
    }
}

class MotorcycleDeliveryFactory extends VehicleFactory {
    @Override
    public Vehicle createVehicle() {
        return new MotorcycleDelivery();
    }
}

class ElectricScooterFactory extends VehicleFactory {
    @Override
    public Vehicle createVehicle() {
        return new ElectricScooter();
    }
}
\end{lstlisting}

\subsubsection{Why This Change?}

Each factory class is a \textbf{Concrete Creator}. These classes carry all of the \texttt{new} keyword usage --- they are the \textit{only} place where a specific vehicle type is instantiated.

\begin{itemize}[leftmargin=*, label=\textcolor{primaryblue}{\textbullet}]
    \item \textbf{Single Responsibility:} Each factory does exactly one thing. \texttt{BusFactory} creates buses. \texttt{TaxiFactory} creates taxis. Nothing more, nothing less.
    \item \textbf{Isolation:} If the \texttt{Bus} class constructor changes (e.g., requires a new parameter), only \texttt{BusFactory} needs updating --- no other part of the system is affected.
    \item \textbf{Extensibility:} To add a \texttt{Train} vehicle, a developer simply adds a \texttt{Train} class and a \texttt{TrainFactory} class. Zero existing code is modified.
\end{itemize}

\subsection{Change 5 --- Client Code with Menu (New Addition)}

\subsubsection{Code Added}

\begin{lstlisting}[caption={Client code --- getFactory() helper method}]
public static VehicleFactory getFactory(int choice) {
    switch (choice) {
        case 1: return new BusFactory();
        case 2: return new TaxiFactory();
        case 3: return new MotorcycleDeliveryFactory();
        case 4: return new ElectricScooterFactory();
        default: return null;
    }
}
\end{lstlisting}

\begin{lstlisting}[caption={Client code --- using the factory (no direct new Vehicle())}]
VehicleFactory factory = getFactory(vehicleChoice);
Vehicle myVehicle = factory.orderVehicle(route);
myVehicle.getVehicleInfo();
double fare = myVehicle.calculateFare(distance);
myVehicle.startTrip();
\end{lstlisting}

\subsubsection{Why This Change?}

The original code had no \texttt{main()} method or client code at all. The refactored version introduces a complete interactive main class that demonstrates the pattern correctly.

The \texttt{getFactory()} method acts as a \textbf{factory selector}. It is the only method that maps a user's input number to a factory class. Critically:

\begin{itemize}[leftmargin=*, label=\textcolor{primaryblue}{\textbullet}]
    \item The \texttt{main()} method \textbf{never calls \texttt{new Bus()} or \texttt{new Taxi()}}. All construction is delegated to factories.
    \item The variable \texttt{myVehicle} is typed as \texttt{Vehicle} (the interface), not as any concrete class. This means the rest of the loop body works identically regardless of which vehicle was selected.
    \item Adding a new vehicle (e.g., \texttt{Train}) only requires: (1) a \texttt{Train} class, (2) a \texttt{TrainFactory} class, and (3) one new \texttt{case} line in \texttt{getFactory()}. The while loop, the fare calculation, the trip start --- all remain untouched.
\end{itemize}

\section{Comparison: Before vs.\ After Refactoring}

\begin{center}
\begin{tabular}{p{5.8cm} p{5.8cm}}
\toprule
\textbf{\color{codered}Before Refactoring} & \textbf{\color{codegreen}After Refactoring} \\
\midrule
No main class or client code & Full interactive main class with menu \\[4pt]
No object creation strategy & Factory Method Pattern applied \\[4pt]
\texttt{new Bus()}, \texttt{new Taxi()} scattered everywhere & All \texttt{new} calls isolated in factory classes \\[4pt]
Adding a new vehicle requires editing main logic & Adding a new vehicle only requires a new class \\[4pt]
No separation of concerns & Creation logic separated from business logic \\[4pt]
Tightly coupled to concrete classes & Client depends only on abstract types \\[4pt]
Difficult to test in isolation & Each factory and vehicle is independently testable \\[4pt]
Violates Open/Closed Principle & Follows Open/Closed Principle \\
\bottomrule
\end{tabular}
\end{center}

\section{Design Principles Applied}

The refactoring enforces three fundamental object-oriented design principles:

\begin{enumerate}[leftmargin=*, label=\textbf{\arabic*.}]
    \item \textbf{Open/Closed Principle (OCP)} \\
    Classes are open for extension (new factories can be added) but closed for modification (existing factories and the abstract class do not need to change).

    \item \textbf{Single Responsibility Principle (SRP)} \\
    Each class has one reason to change. \texttt{Bus} changes only if bus behaviour changes. \texttt{BusFactory} changes only if bus creation changes. \texttt{VehicleFactory} changes only if the shared ordering logic changes.

    \item \textbf{Dependency Inversion Principle (DIP)} \\
    High-level modules (\texttt{main}) depend on abstractions (\texttt{VehicleFactory}, \texttt{Vehicle}), not on concrete implementations (\texttt{Bus}, \texttt{Taxi}).
\end{enumerate}

\section{Conclusion}

The Smart Transportation Management System was successfully refactored from an unstructured prototype into a clean, scalable system using the \textbf{Factory Method Design Pattern}. The original vehicle classes provided a solid foundation through the \texttt{Vehicle} interface, but the system lacked any controlled creation mechanism.

By introducing the abstract \texttt{VehicleFactory} class and four corresponding concrete factory classes, the system now separates object creation from business logic entirely. The client code (the \texttt{main} method) operates only on abstract types, making the system easy to read, extend, and maintain.

The most tangible practical benefit is extensibility: if the transportation company decides to add a new vehicle type --- say, a \texttt{Ferry} or a \texttt{Helicopter} --- a developer needs to write exactly two new classes (\texttt{Ferry} and \texttt{FerryFactory}) and add one case to the factory selector. Not a single line of existing, tested code needs to be changed.

This is the promise of the Factory Method Design Pattern: \textit{grow your system by adding, not by modifying.}

\end{document}